\documentclass[10pt,a4paper]{amsart}
\usepackage[margin=19mm]{geometry}
\usepackage{amsmath,amssymb,mathtools}
\usepackage[scr=rsfs,cal=euler]{mathalfa}
\usepackage{xfrac}
\usepackage[unicode,hidelinks,bookmarks=false]{hyperref}
\newcommand{\ie}{i.e.\ }
\newcommand{\cf}{cf.\ }
\newcommand{\resp}{respectively\ }
\newcommand{\Subsection}{\subsection}
\providecommand{\nobreakdash}{\nobreak\hbox{-}}
\setlength{\emergencystretch}{2em}
\multlinegap0pt
\newtheorem{proofof}{Proofof}
\usepackage{amssymb}
\usepackage{tikz}
\usepackage{mathtools}
\usepackage[shortlabels]{enumitem}
\newtheorem{lemma}{Lemma}
\newcommand{\Bcal}[0]{\ensuremath{{\mathcal B}}}
\newcommand{\Qcal}[0]{\ensuremath{{\mathcal Q}}}
\newcommand{\Zcal}[0]{\ensuremath{{\mathcal Z}}}
\newcommand{\Zed}[0]{\ensuremath{ \mathbb Z}}
\newcommand{\eR}[0]{\ensuremath{ \mathbb R}}
 \newcommand{\eps}{\varepsilon}
\newcommand{\norm}[1]{\ensuremath{\left\|#1\right\|}}
\newcommand{\orig}{\underline{0}}
\title{Thresholds for colouring the random Borsuk graph}
\author{\'Alvaro Acitores Montero, Matthias Irlbeck, Tobias M\"uller, Mat\v{e}j Stehl\'ik}
\date{Source: arXiv:2603.05467v2 (CC BY 4.0)}
\begin{document}
\maketitle

\noindent\emph{Source and licence.} Thresholds for colouring the random Borsuk graph. Version arXiv:2603.05467v2 (CC BY 4.0), distributed by arXiv under the Creative Commons Attribution 4.0 International licence, http://creativecommons.org/licenses/by/4.0/, as recorded in the arXiv licence metadata for this exact version and shown on the abstract page https://arxiv.org/abs/2603.05467v2. Full licence notice: \texttt{licenses/arxiv-2603.05467-CC-BY-4.0.txt}.\par\smallskip

\noindent\textit{Editorial scope.} Counted unit: the article's lemma lem:embed (source0.tex lines 1060--1082). The article's complete original proof of it is delivered with it (source0.tex lines 1672--1794, one \texttt{proof} environment of 123 lines, which the article places away from the statement). No mathematics is written, repaired or re-derived: the statement and the proof are the article's own lines, in the article's own notation, and no retained line is substituted or re-flowed.  Retained uncounted as the source-local context the proof needs: lines 1119--1124; lines 1211--1229.  Nothing was omitted from any retained span, so the package carries no omission record.  No retained line contains a citation, so the wrapper carries no bibliography and no bibliography entry.  No retained line contains a figure environment, an image-inclusion command or drawing code, so no image is delivered and the package declares no asset dependency.  Editorial formatting only, in addition: an AMS \texttt{amsart} preamble that restates the article's own declarations and macros used by the retained lines, namely the environments \texttt{lemma} on separate counters, together with the standard packages those lines need.  The retained environments print in this document as Lemma 1 in order of appearance, whereas the article prints the same items with the numbering of its own full document.  The complete native source of the article as distributed in the arXiv e-print for this exact version is bundled unchanged as \texttt{sources/195805/source0.tex}.
\begin{lemma}\label{lem:boxes}
For every $K, \delta$ there is a $c = c(K,\delta)$ such that, 
setting $\alpha := c \cdot n^{-1/d}, s_0 := \delta \cdot \alpha$, etc., as above, for every $i \geq 2 \log\log n$, 
with probability $1-o_n(1)$, every level-$i$ cube $q \in \Qcal_i$ such that 
$q \subseteq B(\orig,100)$ is good.
\end{lemma}

\begin{lemma}\label{lem:fg}
For every $d\geq 2$ and $k\geq 1$, there exist $t=t(d,k), a_0=a_0(d,k), r_0=r_0(d,k), C=C(d,k) > 0$ such 
that the following holds.
Suppose that $0 < a < a_0, 0 < r < r_0$, the finite set $V \subseteq \eR^d$
and the function $f :S^{d-1} \to (-t,t)$ are such that  
$f$ is odd and $a$-Lipschitz and moreover

$$ \left|B\left( (1+f(u))\cdot u, r \right) \cap V\right| \leq k, $$

\noindent
for all $u \in S^{d-1}$.
Then there also exists a function $g : S^{d-1} \to \eR$ that is odd, $(Ca)$-Lipschitz and satisfies 

$$ |f(u)-g(u)| < Car \quad \text{ and } \quad 
B\left( (1+g(u)) \cdot u, a^{C}\cdot r  \right) \cap V = \emptyset, $$

\noindent 
for all $u \in S^{d-1}$.
\end{lemma}

\begin{lemma}\label{lem:embed}
For every $\eps >0$ there exists a $c=c(\eps)>0$ such that the following holds, 
setting $\alpha = c \cdot n^{-1/d}$.
With probability $1-o_n(1)$, there exists a function 

$$ h : S^{d-1} \to \left(-n^{-0.9/d},n^{-0.9/d}\right),$$ 

\noindent
that is odd, $\eps$-Lipschitz and  

$$ \Zcal \cap B\left( (1+h(u))\cdot u, \eps\cdot\alpha \right) \neq \emptyset,$$ 

\noindent 
for all $u \in S^{d-1}$.
% satisfying the following requirements:
% \begin{enumerate}
%  \item $h(u)+h(-u) = 2$ for all $u \in S^{d-1}$, and; %$\left|h(-u) - \frac{1}{h(u)}\right| \leq \frac{\log^{4}n}{n^{2/d}}$, and;
%  \item $\left|h(u)-h(v)\right| < \eps\cdot\norm{u-v}$ for all $u,v \in S^{d-1}$, and;
%  \item $\Zcal \cap B_{\eR^d}( h(u)\cdot u, \eps\cdot \alpha ) \neq \emptyset$ for all $u \in S^{d-1}$.
%  %, and;
%  %\item $1-\frac{\log^2 n}{n^{1/d}} \leq h(u) \leq 1+\frac{\log^2 n}{n^{1/d}}$ for all $u \in S^{d-1}$.
% \end{enumerate}
\end{lemma}

\begin{proofof}{Lemma~\ref{lem:embed}}
We apply Lemma~\ref{lem:boxes} with parameters $K=K(\eps), \delta=\delta(\eps)$ to be determined 
during the course of the proof. 
For the remainder of the proof, we assume that $\Zcal$ is such that 
the conclusion of Lemma~\ref{lem:boxes} holds. (Which happens with probability $1-o_n(1)$.)
Let $C=C(d,2^d), t=t(d,2^d), a_0=a_0(d,2^d), r_0=r_0(d,2^d)$ be as provided by Lemma~\ref{lem:fg}
applied in dimension $d$ with parameter $k=2^d$.

For $0 \leq j \leq i := \lceil 2\log\log n \rceil$, we set

$$ \Bcal_j := 
\left\{  p \in s_j\Zed^d : \text{$p+[0,s_j]^d$ is bad and contained  in $B(\orig,100)$}\right\}. $$

We construct the sought function $h$ by reverse induction.
Specifically, for $0 \leq j \leq i$, the function $h_j : S^{d-1} \to \eR$ 
will be odd, $(\eps \cdot C^{-j})$-Lipschitz, and satisfy: 

\begin{equation}\label{eq:dem1} 
|h_j(u)| \leq s_{j+1} + s_{j+2} + \dots + s_i, 
\end{equation}

\noindent
and

\begin{equation}\label{eq:dem2} 
B\left( (1+h_j(u))\cdot u, 2\sqrt{d}\cdot s_j \right) \cap \Bcal_j = \emptyset, 
\end{equation}

\noindent
for all $u \in S^{d-1}$.

\noindent
Let us point out that

\begin{equation}\label{eq:sums} 
\begin{array}{rcl}
s_0 + s_1 + \dots + s_i 
& \leq & (i+1)s_i = (i+1) K^{i(i+1)/2} s_0 = e^{O\left( (\log\log n)^2 \right)} \cdot n^{-1/d} \\
& = & 
n^{-1/d+o(1)}. 
\end{array} 
\end{equation}

\noindent
In particular demand~\eqref{eq:dem1} implies that 

$$ h_j(u) \in \left(-n^{-0.9/d},n^{-0.9/d}\right), $$

\noindent
for $n$ sufficiently large. 
Demand~\eqref{eq:dem2} implies that $(1+h_0(u)) \cdot u$ lies in a 
good level-0 cube. In particular 

$$ B( (1+h_0(u))\cdot u, \sqrt{d} \cdot s_0 ) \cap \Zcal \neq \emptyset. $$

\noindent 
Provided $\delta$ was sufficiently small, $\sqrt{d} \cdot s_0 = \sqrt{d} \cdot 
\delta \cdot \alpha < \eps \cdot \alpha$.
So the function $h_0$ can serve as the function $h$ we seek.

As mentioned, we shall establish the existence of $h_0$ by reverse induction. 
We start by defining $h_i$. Clearly $h_i \equiv 0$ does the trick, as $\Bcal_i=\emptyset$.
Now suppose that we've managed to construct $h_{j+1}$ for some $j < i$.

Before proceeding, let us point out that if $h_{j+1}$ satisfies the demands, we have:

\begin{equation}\label{eq:ballsiplus1}
\left| B\left( (1+h_{j+1}(u))\cdot u, \frac{s_{j+1}}{2}\right) \cap \Bcal_j \right| \leq 2^d \quad 
\text{ for all $u \in S^{d-1}$.} 
\end{equation}

\noindent
(Because $B\left( (1+h_{j+1}(u))\cdot u, s_{j+1}/2\right)$ intersects at most $2^d$ level $j+1$ cubes.
None of these can be bad, because that would contradict~\eqref{eq:dem2}.)
We can thus obtain $h_j$ by invoking Lemma~\ref{lem:fg} with 

$$ k=2^d, \quad f=h_{j+1}, \quad V = \Bcal_j, \quad a = \eps\cdot C^{-(j+1)}, \quad 
r = s_{j+1}/2. $$

\noindent
(For $n$ sufficiently large, we are allowed to apply Lemma~\ref{lem:fg} because
$|h_{j+1}| < t$ and $r < r_0$ by~\eqref{eq:sums}, and 
$a < a_0$ assuming without loss of generality $\eps < a_0$.)
We note that 

$$ Ca = \eps \cdot C^{-j}, $$

\noindent
so that $h_j$ is $(\eps C^{-j})$-Lipschitz as required. We also have 

$$ \begin{array}{rcl} 
a^C r 
& = &\left(\eps \cdot C^{-(j+1)}\right)^C \cdot (s_{j+1}/2) \\
& = & \left(\eps \cdot C^{-(j+1)}\right)^C \cdot \frac12 \cdot K^{j+1} \cdot s_{j} \\
& = & \left(\eps^C/2\right) \cdot \left(K/C^C\right)^{j+1} \cdot s_{j} \\
& > & 2\sqrt{d} \cdot s_j,
\end{array} $$

\noindent
where the last line follows because we chose $K$ sufficiently large in the beginning.
So 

$$B\left( (1+h_{j+1}(u))\cdot u, 2\sqrt{d} \cdot s_j\right) \cap \Bcal_j 
\subseteq B\left( (1+h_{j+1}(u))\cdot u, a^C r \right) \cap \Bcal_j 
= \emptyset, $$

\noindent
for all $u \in S^{d-1}$, as
required.
Also note that $\Delta := Car = \eps C^{-j} s_{j+1}/2 < s_{j+1}$, giving

$$ \begin{array}{rcl} 
|h_j(u)| 
& \leq & 
|h_{j+1}(u)| + |h_j(u)-h_{j+1}(u)| \leq 
(s_{j+2} + \dots + s_i) + \Delta \\ 
& < & 
s_{j+1} + s_{j+2} + \dots + s_i. 
\end{array} $$

\noindent
We see that $h_j$ satisfies the demands provided that $h_{j+1}$ does. This concludes the proof.
\end{proofof}

\end{document}
